\documentclass[../main.tex]{subfiles}

\begin{document}

%% ======================================================================
%%  Section 6 – Applications using HDF5
%% ======================================================================
\section{Applications using HDF5}
\label{sec:applications}

Applications treat HDF5 as inert data, but HDF5 files can carry
behavior-relevant metadata: paths, external references, huge
declared shapes, filter/plugin requirements, serialized
application objects, and model-layer definitions. That creates
integration bugs even when the HDF5 library is behaving as
designed.

The HDF Group's security policy~\cite{hdf5-sec-pol} already attempts the right scope
distinction: core libhdf5, official tools, Remote Code Execution (RCE)
via parsing/API abuse, and supply-chain compromise are in scope;
third-party plugins and application-level usage are usually
triaged or referred to their maintainers.

\begingroup
\small
\setlength{\tabcolsep}{3pt}
\renewcommand{\arraystretch}{1.12}
\begin{longtable}{@{}p{0.21\linewidth}p{0.22\linewidth}p{0.51\linewidth}@{}}
\caption{Representative application-level HDF5 mechanisms}
\label{tab:application-hdf5-failure-families}\\
\hline
\textbf{Mechanism} & \textbf{Consequence} & \textbf{Boundary evidence} \\
\hline
\endfirsthead
\hline
\textbf{Mechanism} & \textbf{Consequence} & \textbf{Boundary evidence} \\
\hline
\endhead
\hline
\endfoot
Malformed metadata in C paths & Memory faults or termination &
Recurring core/tool records; the bounded official-tool example is
\code{CORE-TOOL-001}~\cite{cve-2026-34734,hdf5-ghsa-w7v2-9cmr-pwwj}. \\
Application deserialization & Code execution &
Legacy Keras HDF5 loading bypassed \code{safe_mode}; application objects, not
the format, supplied execution~\cite{cve-2025-9905,keras-ghsa-36rr-ww3j-vrjv}. \\
External resources & File disclosure or traversal &
File-selected paths were resolved with application authority in Keras and the
Hugging Face service~\cite{cve-2026-1669,keras-ghsa-3m4q-jmj6-r34q,
huggingface-agent-intrusion-2026}. \\
Declared size or shape & Memory exhaustion &
Keras allocated from an extreme HDF5-declared shape before enforcing a budget
~\cite{cve-2026-0897,keras-ghsa-mgx6-5cf9-rr43}. \\
Dynamic plugin search & Unintended code loading &
TensorFlow searched an unsecured plugin location
~\cite{cve-2026-2492,zdi-26-116,tensorflow-commit-46e7f7f}. \\
Vendored/stale copies & Inherited defects and ambiguous applicability &
Downstreams may statically link or lag HDF5, as illustrated by matio and an
Apple removal~\cite{cve-2021-36977,cve-2021-31009}. \\
\end{longtable}
\endgroup

\iffalse
\begingroup
\scriptsize
\setlength{\tabcolsep}{2pt}
\renewcommand{\arraystretch}{1.15}
\setlength{\LTleft}{-0.7in}
\setlength{\LTright}{-0.7in}

\begin{longtable}{@{}p{0.18\widetablewidth}p{0.18\widetablewidth}p{0.32\widetablewidth}p{0.28\widetablewidth}@{}}
\caption{Representative application-level HDF5 failure mechanisms}
\label{tab:application-hdf5-failure-families}\\
\hline
\textbf{Family} & \textbf{Typical failure} & \textbf{Examples} & \textbf{Application-boundary mechanism} \\
\hline
\endfirsthead

\caption[]{Representative application-level HDF5 failure mechanisms, continued}\\
\hline
\textbf{Family} & \textbf{Typical failure} & \textbf{Examples} & \textbf{Application-boundary mechanism} \\
\hline
\endhead

\hline
\endfoot

\hline
\endlastfoot

malformed file metadata reaches C memory code &
buffer overflow, over-read, UAF, double-free, null dereference, divide-by-zero &
many HDFGroup records, including 2024 and 2025 issues in object headers, datatypes, dataspaces, filters, heaps, free-space metadata, and cache images &
HDF5 is a complex binary metadata graph, not a flat byte container\\
application deserialization layered on HDF5 &
arbitrary code execution &
Keras CVE-2025-9905: crafted \code{.h5/.hdf5} model archives can execute code
because legacy \code{.h5} loading did not honor \code{safe_mode=True} for
\code{Lambda}-layer pickled code~\cite{cve-2025-9905,keras-ghsa-36rr-ww3j-vrjv}. &
ML frameworks often use HDF5 as a model container, then deserialize
framework-level objects from it. \\
external references and file path indirection &
arbitrary local file read, path traversal, data exfiltration &
Keras CVE-2026-1669: crafted \code{.keras} model file using HDF5 external
dataset references could read local files~\cite{cve-2026-1669,keras-ghsa-3m4q-jmj6-r34q}. &
HDF5 supports links and external storage mechanisms. External links store both
target filename and object path; external storage can resolve raw-data files
through prefixes and current working directory rules~\cite{hdf5-data-model}. \\
unbounded resource allocation &
memory exhaustion / interpreter crash &
Keras CVE-2026-0897: crafted \code{.keras} archive with
\code{model.weights.h5} declaring an extremely large dataset shape could exhaust
memory and crash Python~\cite{cve-2026-0897,keras-ghsa-mgx6-5cf9-rr43}. &
HDF5 metadata can declare logical shapes much larger than the actual useful
payload; application code may allocate based on declared shape before policy
checks. \\
dynamic plugin loading and search paths &
local privilege escalation / arbitrary code loading &
TensorFlow CVE-2026-2492: HDF5 plugin handling loaded plugins from an unsecured
location, enabling local privilege escalation
\cite{cve-2026-2492,zdi-26-116,tensorflow-commit-46e7f7f}. &
HDF5 supports dynamic plugin loading, and plugin paths can be affected by
defaults or \code{HDF5_PLUGIN_PATH}; dynamic loading can also be disabled via
\code{H5PLset_loading_state()} or
\code{HDF5_PLUGIN_PRELOAD=::}~\cite{hdf5-dynamic-plugins}. \\
vulnerable or stale embedded HDF5 copies &
inherited CVEs, unpatched parser behavior &
matio CVE-2021-36977 involved a heap overflow related to use of HDF5
1.12.0; Apple CVE-2021-31009 says multiple issues were addressed by removing
HDF5~\cite{cve-2021-36977,cve-2021-31009}. &
Downstream products vendor, statically link, or lag HDF5 releases;
vulnerability scanners then attribute HDF5 defects to the application. \\
helper-tool exposure &
crashes or code execution through inspection/conversion utilities &
\code{h5dump} appears in older and newer issues; CVE-2026-34734 describes a
heap-use-after-free in the \code{h5dump} helper utility triggered by a malicious
\code{.h5} file~\cite{cve-2026-34734,hdf5-ghsa-w7v2-9cmr-pwwj}. &
Applications often call \code{h5dump}, \code{h5repack}, or converters as
``safe inspection'' tools on untrusted files. \\
\end{longtable}
\endgroup
\fi

\subsection{Boundary Case Study: External Raw Storage and the Hugging Face Incident}
\label{sec:huggingface-external-storage-case}

The July 2026 Hugging Face incident is a concrete example of an application
trust-boundary failure involving an intentional HDF5 capability. Hugging Face
reports that its dataset service processed attacker-supplied, valid HDF5 files
whose datasets declared local paths as external raw storage. The worker read
\code{/proc/self/environ} and source files with its own filesystem permissions
and returned the bytes as dataset rows. Hugging Face expressly distinguishes
this file-disclosure vector, which executed no code, from a separate Jinja2
injection that supplied code execution~\cite{huggingface-agent-intrusion-2026}.

The HDF Group likewise concludes that external raw storage was the mechanism,
not an HDF5 parser vulnerability: the documented feature behaved as designed,
while the authority-bearing service acted as a confused deputy. Updating
libhdf5 alone would not remove that service boundary; the relevant controls are
fail-closed external-resource policy, preflight inspection, path confinement,
least-privileged workers, and credential isolation
\cite{hdfgroup-mechanism-not-vulnerability-2026}.

The h5py community expressly considered the same boundary. Maintainers described
per-dataset external-storage checks as possible but laborious, recommended
sandboxing genuinely untrusted files, opened upstream issue \#6618 because
h5py had no direct enforcement mechanism, and discussed a simpler detection
API. Issue \#6618 in turn asks for unified controls at or before external-resource
resolution across raw storage, external links, and virtual-dataset sources. The
discussion records policy and API proposals, including use of the external
\code{h5policy} prototype as a proving ground; this report does not treat those
proposals as delivered controls~\cite{h5py-gh-issue-2891,hdf5-gh-issue-6618}.

Table~\ref{tab:huggingface-authority-chain} separates the observed incident
record from the control inference. That separation matters because the file was
not reported to exploit memory corruption, and the service's authority --- not
the input's authority --- made the disclosure possible.

\begingroup
\small
\setlength{\tabcolsep}{3pt}
\renewcommand{\arraystretch}{1.15}
\begin{longtable}{@{}p{0.17\linewidth}p{0.47\linewidth}p{0.30\linewidth}@{}}
\caption{Authority and control chain in the Hugging Face incident}
\label{tab:huggingface-authority-chain}\\
\hline
\textbf{Stage} & \textbf{Evidence and boundary} & \textbf{Audit consequence} \\
\hline
\endfirsthead
\caption[]{Authority and control chain in the Hugging Face incident, continued}\\
\hline
\textbf{Stage} & \textbf{Evidence and boundary} & \textbf{Audit consequence} \\
\hline
\endhead
\hline
\endfoot

Resource selection &
The remote party supplied a valid HDF5 file whose external raw-storage metadata
named local paths. The originator selected the resource but did not bring the
worker's filesystem authority. &
Validity and parser correctness do not answer whether the application may honor
a file-selected resource. External access needs an explicit policy decision. \\

Authority-bearing resolution &
The dataset worker resolved the paths with its own permissions and read
\code{/proc/self/environ} and source files. The intentional HDF5 mechanism
carried the request across the service boundary. &
Deny or constrain resolution before the read; independently reduce worker
authority, isolate credentials, and confine filesystem access. \\

Observable consequence &
The service returned the bytes as dataset rows. Hugging Face distinguishes this
file disclosure from the separate Jinja2 injection that supplied code
execution~\cite{huggingface-agent-intrusion-2026}. &
Rate the demonstrated disclosure, not an unreported HDF5-code-execution path.
Constrain output and retain resolution logs so attempted access is observable. \\

Reusable-control gap &
The affected renderer was changed, while h5py and HDF5 discussions sought a
less laborious detection or enforcement surface spanning raw storage, external
links, and virtual-dataset sources~\cite{h5py-gh-issue-2891,hdf5-gh-issue-6618}. &
Service containment and reusable library support are different closure claims.
Until enforcement ships and is integrated, sandboxing and application policy
remain necessary. \\

Closure evidence &
A preflight result can become stale if the file, path environment, or opened
object changes before use. A detector alone also cannot decide which worker
authority is acceptable. &
Bind the decision to the same handle or immutable file identity, recheck at
resolution, and test denial across every enabled external-resource family. \\

\end{longtable}
\endgroup

Formal finding \code{APP-SEC-001} in
Appendix~\ref{app:findings-register} rates the
incident at the affected service boundary. The similar Keras external-dataset
failure is separately registered as \code{APP-SEC-003}; recurrence of a
mechanism does not erase the differences in product, version, and deployment
scope.

\paragraph{The unifying root cause.}
The repeated mistake is \emph{trust-boundary confusion}: treating HDF5 as inert
data even though it can influence allocation, shape, traversal, codec/plugin
loading, external-file lookup, conversion, object reconstruction, and tool
output. The SSP SIG separately names the same family as I9 in its
independent-implementation model and W8 in its Python-wrapper model
~\cite{hdf5-indie-model,hdf5-python-wrapper-model}. That is internal
consistency between programme artifacts, not independent validation. The
failure occurs when a model loader, ingestion service, workflow, viewer, or
plugin resolver interprets file-controlled metadata with authority the file's
originator did not possess.

\subsection{Audit Findings and Mitigation Planning}

The formal application findings are \code{APP-SEC-001} through
\code{APP-SEC-005} in Appendix~\ref{app:findings-register}. Their ratings apply
to the bounded versions and deployments stated there, not to Keras,
TensorFlow, Hugging Face, or applications using HDF5 as undifferentiated
projects.

\paragraph{Priority reconciliation.}\mbox{}\\
Any equivalent deployment that still gives untrusted HDF5 input automatic
access to file-selected external resources follows the Immediate route for
\code{APP-SEC-001}: contain it before the next job and no later than 24 hours
after identification. Confirmed affected scope for \code{APP-SEC-002},
\code{APP-SEC-003}, and \code{APP-SEC-005} receives Near-term treatment;
\code{APP-SEC-004} receives Planned treatment. The historical Hugging Face
rating is not a current residual-risk rating of that service. Application
patches and compensating controls in Phase~1 must not wait for the reusable
core APIs and integrations proposed for Phase~2.

Mitigation has to be layered. Patching \texttt{libhdf5}
alone will not fix Keras model deserialization, TensorFlow plugin
path handling, or application-level external-reference policy.

The \code{APP-S} identifiers below identify recommendations in the downstream-
application track. Their numbers do not express risk, urgency, effort,
dependency, or implementation order.

\begingroup
\small
\setlength{\tabcolsep}{3pt}
\renewcommand{\arraystretch}{1.1}
\begin{longtable}{@{}p{0.12\linewidth}p{0.26\linewidth}p{0.56\linewidth}@{}}
\caption{Application-boundary control set}
\label{tab:downstream-hdf5-safe-policies}\\
\hline
\textbf{ID} & \textbf{Outcome} & \textbf{Minimum control} \\
\hline
\endfirsthead
\hline
\textbf{ID} & \textbf{Outcome} & \textbf{Minimum control} \\
\hline
\endhead
\hline
\endfoot

\code{APP-S1} & Untrusted-HDF5 profile &
Default-deny dynamic plugins and external resources; cap rank, objects,
metadata, traversal, logical bytes, allocation, and filter expansion. \\
\code{APP-S2} & Enforcement-coupled preflight &
Return machine-readable reasons before allocation or deserialization, bind the
decision to the same handle or immutable file identity, and recheck targets at
runtime. Existing validator prior art does not supply that coupling. \\
\code{APP-S3} & Safe ML-loader defaults &
Reject executable legacy reconstruction, external datasets, dynamic plugins,
and allocation beyond aggregate and per-object budgets. \\
\code{APP-S4} & External resources opt-in &
Reject absolute paths, traversal, and working-directory lookup; require an
explicit sandbox root or allowlist and log each resolution. \\
\code{APP-S5} & Plugins governed as code &
Set loading policy explicitly; use protected allowlisted directories and
provenance; isolate and minimize privilege where a plugin is indispensable. \\
\code{APP-S6} & Downstream provenance &
Report HDF5 version and linkage, enabled filters/connectors, plugin-path and
external-resource policy, and loader safe-mode behavior. \\
\code{APP-S7} & Advisory scope tags &
Distinguish core parser, official or legacy tool, plugin loading, external
resource, downstream deserialization/resource policy, vendored library,
non-HDF5, and duplicate/rejected records. Tags support triage but do not rate
risk. \\
\end{longtable}
\endgroup

\paragraph{Preflight must remain bound to enforcement.}
A metadata scan is advisory unless its decision is attached to the same open
handle or a verified immutable file identity and the runtime fails closed when
the file, policy, or external-resource target changes. The validator should
report reasons such as an out-of-sandbox path, an excessive logical size, a
required plugin, or executable application metadata before the loader allocates
arrays or constructs application objects. Runtime hooks must then enforce the
same limits at external resolution, traversal, filter activation, and object
construction.

\paragraph{Use current controls while shared APIs mature.}
Applications can already disable dynamic plugins with
\code{H5PLset_loading_state(0)} or \code{HDF5_PLUGIN_PRELOAD=::}; service
launchers can clear or constrain \code{HDF5_PLUGIN_PATH} and
\code{HDF5_EXTFILE_PREFIX}; and process isolation can remove filesystem,
credential, and network authority. Those controls should be published as
deployment recipes now. The Phase~2 profile and external-resource APIs make the
policy cohesive and binding-friendly; they do not justify waiting to contain an
identified deployment.

\paragraph{Deployment minimum.}
Unknown-file viewers and upload services require sandboxing, external-resource
denial, plugin denial, and metadata/output budgets. Model loaders additionally
reject executable reconstruction and enforce aggregate tensor budgets. HPC and
desktop workflows separate trusted internal files from imports, pin protected
plugin directories, and make any external-resource opt-in explicit. Bindings
must expose these controls through APIs rather than environment variables alone.

Viewers, model loaders, upload services, HPC workflows, desktop applications,
bindings, and distributors apply this set differently, but the common rule is
the same: isolate unknown files, disable unneeded authority, enforce resource
budgets before allocation, and publish the exact HDF5 boundary in use. The
authoritative owner, work package, effort, and acceptance evidence are in
Appendices~\ref{app:recommendation-register} and
\ref{app:phased-roadmap}.

\iffalse
\subsubsection*{Recommendation APP-S1: Define an “untrusted HDF5” profile}

HDF5 should offer, and downstreams should use, a strict profile
for any file received from users, networks, archives, model hubs,
object stores, email, web uploads, or third-party datasets.

Default behavior for that profile:
\begin{verbatim}
dynamic plugins disabled
external links disabled or sandboxed
external raw dataset storage disabled or sandboxed
VDS source traversal restricted
maximum rank enforced
maximum element count enforced
maximum allocation size enforced
maximum metadata bytes enforced
maximum chunk count enforced
maximum filter expansion ratio enforced
cycles and deep graph traversal rejected
unknown or legacy-dangerous messages rejected
    unless explicitly allowed    
\end{verbatim}

Today, pieces exist: plugin loading can be controlled, and all dynamic plugins
can be disabled with \code{H5PLset_loading_state(0)} or \code{HDF5_PLUGIN_PRELOAD=::}.
But this needs to become a cohesive security profile, not a set of hidden knobs.

\subsubsection*{Recommendation APP-S2: Ship a preflight validator}

\noindent\emph{Prior art:} a working preflight validator already exists; see
Table~\ref{tab:core-prior-art-rescope}. The open work is the enforcement
coupling described below, not the validator itself.

Create a validator that reads metadata before the application allocates large
arrays or deserializes model objects. Preflight is not a stand-alone
authorization: its result must be bound to the same open handle or verified
immutable file identity used by runtime loading, and runtime hooks must fail
closed if the file, policy, or external-resource target changes.

Example output:
\begin{minted}{json}
{
  "eligible_under_profile": false,
  "reasons": [
    "external dataset reference outside sandbox",
    "dataset declares 8.7 TB logical size",
    "requires dynamic filter plugin",
    "contains legacy Keras Lambda layer metadata"
  ]
}    
\end{minted}

For core HDF5, this validator should reason about file-format invariants.
For downstreams, it should expose application hooks:

\begin{verbatim}
on_external_reference(path)
on_dataset_shape(path, dtype, shape, logical_bytes)
on_filter_required(filter_id)
on_link_traversal(source, target)
on_attribute(name, size)
\end{verbatim}

Keras and TensorFlow could then enforce model-specific policy before creating Python objects.

\subsubsection*{Recommendation APP-S3: Disable dangerous HDF5 features by default in ML model loading}

ML loaders should not behave like general HDF5 viewers.

For Keras/TensorFlow-style use:
\begin{verbatim}
do not allow Lambda or pickled code from legacy .h5
do not allow external dataset references
do not allow dynamic HDF5 plugins
do not eagerly allocate declared tensors before checking total
    byte budget
do not resolve paths relative to current working
    directory
do not honor HDF5_PLUGIN_PATH in model loading
    contexts
\end{verbatim}   
These loader defaults directly address \code{APP-SEC-002} through
\code{APP-SEC-004}. The plugin clauses provide defense in depth for
\code{APP-SEC-005}; its primary linked controls are \code{APP-S5} and
\code{CORE-S10}.

\subsubsection*{Recommendation APP-S4: Make external references opt-in}

External links and external raw dataset storage are legitimate HDF5 features,
but they are unsafe defaults for hostile content.
HDF5 documentation states that external links store a target
filename and object path, and external dataset storage search
behavior can involve the current working directory, absolute
paths, relative paths, \code{${ORIGIN}}, and \code{HDF5_EXTFILE_PREFIX}.

Mitigation:
\begin{verbatim}
reject absolute external paths
reject path traversal
reject current-working-directory lookup
require explicit sandbox root
require explicit allowlist by extension or path
log every external resolution
never follow external references during metadata
    preview
\end{verbatim}

For services, the safest default is: \code{external references = off}.

\subsubsection*{Recommendation APP-S5: Treat plugins as executable code}

HDF5 filter plugins, VOL connectors, and VFDs are not data.
They are code. The HDF5 docs state that plugins are shared
libraries and that the default path can be overwritten by
\code{HDF5_PLUGIN_PATH}. That is enough to create an
application-level vulnerability if a privileged process uses an
attacker-controlled path.

Mitigation:
\begin{verbatim}
clear HDF5_PLUGIN_PATH, HDF5_PLUGIN_PRELOAD,
    HDF5_EXTFILE_PREFIX in services
set plugin loading state explicitly at process start
allow plugins only from root/admin-owned directories
reject writable plugin directories
require hashes or signatures for plugins
never load plugins in setuid, service, notebook-hub,
    web-worker, or privileged contexts unless explicitly
    allowlisted
\end{verbatim}

Where a required plugin cannot be disabled, run the file-processing boundary
with least privilege and process isolation in addition to enforcing protected
paths, explicit allowlisting, and resource limits.

\subsubsection*{Recommendation APP-S6: Require downstream SBOM and version reporting}

Downstream projects should report:
\begin{verbatim}
HDF5 version
linkage type: static / dynamic / vendored / system
enabled filters
enabled VOL/VFD connectors
plugin search path policy
external-reference policy
safe-mode policy
\end{verbatim}

This reduces a practical applicability ambiguity: a user seeing “HDF5 CVE” in a
scanner often cannot tell whether they are affected through
\code{libhdf5}, an official tool, a vendored copy, a Python wheel,
a model loader, or a plugin. Reporting supplies census inputs; it does not by
itself establish that a deployment is affected, compatible, or remediated.

\subsubsection*{Recommendation APP-S7: Split CVE classification into scope tags}

For HDF5-related CVE triage, use explicit tags:
\begin{verbatim}
HDF5-core-parser
HDF5-official-tool
HDF5-legacy-tool
HDF5-plugin-loading
HDF5-external-reference
HDF5-downstream-deserialization
HDF5-downstream-resource-policy
HDF5-vendored-library
not-HDF5
duplicate/rejected
\end{verbatim}

This prevents scope collapse and supplies an input to triage. Each bounded
scenario must still be rated and routed under the report-wide rubric; a tag
does not assign priority. A heap overflow in
\code{H5T__conv_struct_opt} is not the same class as Keras
ignoring \code{safe_mode} for legacy \code{.h5} files, even though both
mention HDF5.

\subsubsection*{A practical mitigation matrix for downstream maintainers}

\begingroup
\small
\setlength{\tabcolsep}{3pt}
\renewcommand{\arraystretch}{1.15}
\setlength{\LTleft}{-0.7in}
\setlength{\LTright}{-0.7in}
\begin{longtable}{@{}p{0.25\widetablewidth}p{0.71\widetablewidth}@{}}
\caption{Baseline protective controls requiring deployment-specific validation}
\label{tab:downstream-hdf5-safe-policies}\\
\hline
\textbf{Downstream use case} & \textbf{Baseline protective controls} \\
\hline
\endfirsthead

\caption[]{Baseline protective controls requiring deployment-specific validation, continued}\\
\hline
\textbf{Downstream use case} & \textbf{Baseline protective controls} \\
\hline
\endhead

\hline
\endfoot

\hline
\endlastfoot

viewing unknown \code{.h5} files &
run in a sandbox; disable plugins; disable external references; cap metadata
and output size \\
loading ML models &
use a verified fixed loader or reject or contain affected legacy formats;
disable executable reconstruction, external references, and plugins; enforce
aggregate and per-object budgets before allocation \\
cloud upload service &
enforcement-coupled preflight and runtime denial bound to the same file
identity; no plugin loading or external file access; resource limits and
process isolation \\
HPC workflow &
separate trusted internal files from imported files; scrub environment; pin
plugin directories; use module-level HDF5 version floors \\
desktop scientific app &
default-deny external references; permit only explicit, path-confined opt-in
after validation; ship applicable fixes; avoid writable global plugin dirs \\
package/distribution &
publish SBOM; expose HDF5 version; confirm applicability and meet the
finding-specific response horizon, retaining tested containment when a rebuild
is unavailable \\
library binding &
expose security controls in the binding API, not just through environment
variables \\
\end{longtable}
\endgroup

\paragraph{Reusable controls The HDF Group can provide to reduce downstream exposure:}
For reusable downstream prevention, parser hardening should be paired with
enforceable controls that applications can use easily. Publish recipes using
controls available today in parallel, and update them as the Phase~2 APIs ship.
This delivery sequence does not replace or defer Immediate or Near-term
application treatment.

Recommended deliverables:
\begin{verbatim}
H5Pset_security_profile(fapl, H5_SECURITY_PROFILE_UNTRUSTED)
H5Pset_external_reference_policy(fapl, H5_EXTREF_DENY
    | H5_EXTREF_SANDBOX)
H5Pset_plugin_policy(fapl, H5_PLUGIN_DENY
    | H5_PLUGIN_ALLOWLIST)
H5Pset_max_logical_dataset_bytes(fapl, n)
H5Pset_max_metadata_bytes(fapl, n)
H5Pset_max_object_count(fapl, n)
H5Pset_max_link_traversal_depth(fapl, n)
h5validate --profile untrusted --json file.h5
\end{verbatim}

Maintain downstream recipes for Keras, TensorFlow, h5py, netCDF, MATLAB, R,
Java, and cloud services as current controls and standardized APIs evolve.
\fi

\end{document}
